\section{Quantifiers: The Game of For Every and There Exists}
\chapterline

In Chapter~3, we saw that variables in mathematics serve distinct logical roles: an \emph{arbitrary} variable represents any generic element of a domain, while a \emph{chosen} variable is a specific witness constructed to establish an existential claim.

Yet in mathematical statements, variables do not appear only in isolation. They appear in pairs, triplets, and chains. When multiple variables interact in the same claim, who sets the terms first? Which variable choices may depend on earlier ones?

The answer is governed by two symbols called \textbf{quantifiers}---operators that specify \emph{how many} elements in a collection satisfy a given property:
\[
\forall \quad \text{(``for every'' / ``for all'')}, \qquad\qquad \exists \quad \text{(``there exists'' / ``for at least one'')}.
\]
The order in which these symbols appear determines the logical structure of a proof. Reversing two quantifiers is not a minor syntactic detail; it frequently transforms a true theorem into a false statement.

\begin{historybox}
In the early nineteenth century, Augustin-Louis Cauchy was striving to put calculus on solid ground. Working in 1821 with informal notions of limits, Cauchy stated in his \emph{Cours d'Analyse} that the sum of any convergent series of continuous functions must be continuous. However, Joseph Fourier had already observed that trigonometric series of sines and cosines could converge to discontinuous functions, such as square waves. For decades, mathematicians debated how a sum of continuous terms could yield a discontinuous limit.

The contradiction was resolved in 1847 when Philipp Ludwig von Seidel and George Gabriel Stokes independently realized that Cauchy had conflated two different orders of quantifiers. If you pick any single point on the horizontal axis, the sequence of functions may settle down at that specific point ($\forall x \, \exists N$, known as \emph{pointwise convergence}). That does not guarantee that the sequence settles down at the same rate across the entire domain ($\exists N \, \forall x$, formalized by Karl Weierstrass as \emph{uniform convergence}).

When Gottlob Frege (1879) and Charles Sanders Peirce (1885) invented explicit symbolic notation for quantifiers, this linguistic ambiguity vanished. Quantifiers gave mathematicians the precision to distinguish pointwise convergence from uniform convergence. Alternating quantifiers remain central in modern logic and theoretical computer science: the \emph{Polynomial Hierarchy} classifies computational problems by how many alternating $\forall$ and $\exists$ decisions they require. Whether adding more quantifier alternations strictly expands the class of solvable problems, or whether the hierarchy collapses to a finite level, remains a fundamental open problem in complexity theory.
\end{historybox}

To see why quantifier order is essential, it helps to view quantified statements through game-theoretic semantics, where truth is evaluated as an interaction between two participants.

\subsection{Quantifiers as a two-player game}

Imagine a game played on a mathematical domain between two players:
\begin{itemize}
\item \textbf{The Universal Challenger ($\forall$)}: A challenger seeking a counterexample to the claim. Whenever a statement specifies $\forall x$, the Challenger selects an element $x$ from the domain in an attempt to make the condition fail.
\item \textbf{The Existential Prover ($\exists$)}: The prover attempting to verify the claim. Whenever a statement specifies $\exists y$, the Prover must produce an element $y$ that satisfies the required condition.
\end{itemize}

The fundamental rule of this game is simple: \textbf{the order of the quantifiers dictates the order of play}. Whoever appears first moves first. Whoever moves second plays with full knowledge of the first player's move.

\subsection{Pointwise dependence versus uniform choice}

Let us compare these two roles on the real numbers $\R$ with the relation $x + y = 0$.

\textbf{Game 1}: $\forall x \in \R, \, \exists y \in \R \text{ such that } x + y = 0$.
\begin{itemize}
\item The Challenger plays first, choosing any real number $x$.
\item You play second. Given their choice of $x$, you respond with $y = -x$.
\item Check the condition: $x + (-x) = 0$. Because you can produce a valid $y$ for any $x$ the Challenger selects, the statement is \textbf{True}.
\end{itemize}

\textbf{Game 2}: $\exists y \in \R \text{ such that } \forall x \in \R, \, x + y = 0$.
\begin{itemize}
\item You must play first: you must choose a single real number $y$ before the Challenger reveals their move.
\item Suppose you choose $y = 0$.
\item The Challenger now plays second, inspecting your $y$. Knowing your choice is $0$, they pick $x = 5$.
\item Check the condition: $5 + 0 = 5 \neq 0$.
\item Can any choice of $y$ work? No: for any fixed $y$, the Challenger simply picks $x = 1 - y$, forcing $x + y = 1 \neq 0$. The statement is \textbf{False}.
\end{itemize}

In Game 1, the choice of $y$ depends on $x$: your response can vary with each input. In Game 2, you must select a single uniform value of $y$ that works for every possible $x$.

\begin{recoverybox}
\textbf{Mental check:} What happens if we change the domain of Game~1 from the real numbers $\R$ to the natural numbers $\N = \{1, 2, 3, \ldots\}$?
\[
\forall x \in \N, \, \exists y \in \N \text{ such that } x + y = 0.
\]
Can the Prover still win?
\emph{Answer:} No. The Challenger picks $x = 1$. The only real number satisfying $1 + y = 0$ is $y = -1$, but $-1 \notin \N$. No such natural number exists, so the statement is False. The domain of a quantifier limits the allowable choices for each variable.
\end{recoverybox}

\subsection{Visualizing the game on a grid}

We can make this distinction visible on a grid where rows represent inputs chosen by the Challenger ($x$) and columns represent responses chosen by the Prover ($y$). A cell $(x, y)$ is lit up whenever the predicate $P(x, y)$---a mathematical condition whose truth depends on the pair $(x, y)$---is satisfied.

\begin{center}
\begin{tikzpicture}[
    font=\small,
    cell/.style={draw=gray!30, fill=white, minimum size=0.7cm, inner sep=0pt},
    hitpoint/.style={circle, fill=targetgreen, inner sep=2.2pt},
    mastercol/.style={fill=sourceblue!15, draw=sourceblue!60, thick}
]

% Panel A: Forall x Exists y
\begin{scope}[xshift=0cm]
    \node[anchor=west, font=\footnotesize\scshape, text=targetgreen!80!black] at (-0.8, 4.7) {$\forall x \, \exists y \, P(x,y)$: Pointwise Response};
    
    % Axis role labels
    \node[font=\scriptsize\bfseries, text=gray!60, rotate=90] at (-0.3, 2.0) {Challenger picks $x$};
    \node[font=\scriptsize\bfseries, text=gray!60] at (2.0, -0.2) {Prover responds with $y$};

    % Grid background
    \foreach \r in {1,2,3,4} {
        \foreach \c in {1,2,3,4} {
            \node[cell] at (\c*0.8, \r*0.8) {};
        }
    }
    
    % Axis labels
    \node[font=\scriptsize, text=gray!80] at (0.2, 3.2) {$x_4$};
    \node[font=\scriptsize, text=gray!80] at (0.2, 2.4) {$x_3$};
    \node[font=\scriptsize, text=gray!80] at (0.2, 1.6) {$x_2$};
    \node[font=\scriptsize, text=gray!80] at (0.2, 0.8) {$x_1$};
    
    \node[font=\scriptsize, text=gray!80] at (0.8, 0.2) {$y_1$};
    \node[font=\scriptsize, text=gray!80] at (1.6, 0.2) {$y_2$};
    \node[font=\scriptsize, text=gray!80] at (2.4, 0.2) {$y_3$};
    \node[font=\scriptsize, text=gray!80] at (3.2, 0.2) {$y_4$};
    
    % Green hits: at least one per row
    \node[hitpoint] at (1.6, 0.8) {}; % (x1, y2)
    \node[hitpoint] at (3.2, 1.6) {}; % (x2, y4)
    \node[hitpoint] at (0.8, 2.4) {}; % (x3, y1)
    \node[hitpoint] at (2.4, 3.2) {}; % (x4, y3)
    
    % Row indicators: green checkmarks
    \foreach \r in {0.8, 1.6, 2.4, 3.2} {
        \node[font=\footnotesize, text=targetgreen!80!black] at (3.9, \r) {$\checkmark$};
    }
    
    \node[text width=4.6cm, align=left, font=\scriptsize, text=gray!80!black] at (2.0, -0.9) {
        \textbf{Requirement:} Every single row must contain \emph{at least one} solution point. The column can vary from row to row.
    };
\end{scope}

% Divider line
\draw[draw=gray!30, dashed, thick] (5.3, -1.4) -- (5.3, 4.8);

% Panel B: Exists y Forall x
\begin{scope}[xshift=6.6cm]
    \node[anchor=west, font=\footnotesize\scshape, text=sourceblue!80!black] at (-0.8, 4.7) {$\exists y \, \forall x \, P(x,y)$: Uniform Master Key};
    
    % Axis role labels
    \node[font=\scriptsize\bfseries, text=gray!60, rotate=90] at (-0.3, 2.0) {Challenger tests with $x$};
    \node[font=\scriptsize\bfseries, text=sourceblue!80!black] at (2.0, -0.2) {Prover picks $y^*$ first};

    % Highlighted master column y3
    \draw[mastercol, rounded corners=2pt] (2.05, 0.45) rectangle (2.75, 3.55);
    
    % Grid background
    \foreach \r in {1,2,3,4} {
        \foreach \c in {1,2,3,4} {
            \node[cell] at (\c*0.8, \r*0.8) {};
        }
    }
    
    % Axis labels
    \node[font=\scriptsize, text=gray!80] at (0.2, 3.2) {$x_4$};
    \node[font=\scriptsize, text=gray!80] at (0.2, 2.4) {$x_3$};
    \node[font=\scriptsize, text=gray!80] at (0.2, 1.6) {$x_2$};
    \node[font=\scriptsize, text=gray!80] at (0.2, 0.8) {$x_1$};
    
    \node[font=\scriptsize, text=gray!80] at (0.8, 0.2) {$y_1$};
    \node[font=\scriptsize, text=gray!80] at (1.6, 0.2) {$y_2$};
    \node[font=\scriptsize\bfseries, text=sourceblue!90!black] at (2.4, 0.2) {$y^*$};
    \node[font=\scriptsize, text=gray!80] at (3.2, 0.2) {$y_4$};
    
    % Green hits all down column 3
    \node[hitpoint] at (2.4, 0.8) {};
    \node[hitpoint] at (2.4, 1.6) {};
    \node[hitpoint] at (2.4, 2.4) {};
    \node[hitpoint] at (2.4, 3.2) {};
    
    % Column indicator
    \draw[-{Latex[length=2mm]}, very thick, sourceblue] (2.4, 4.0) -- (2.4, 3.65);
    \node[above=1pt, font=\scriptsize\bfseries, text=sourceblue!90!black] at (2.4, 4.0) {Single key $y^*$};
    
    \node[text width=4.6cm, align=left, font=\scriptsize, text=gray!80!black] at (2.0, -0.9) {
        \textbf{Requirement:} There must exist \emph{one specific column} $y^*$ that is completely filled from top to bottom.
    };
\end{scope}

\end{tikzpicture}
\end{center}

Notice the visual distinction:
\begin{itemize}
\item To satisfy $\forall x \, \exists y \, P(x, y)$, at least one point in each row must be satisfied, even if the column varies from row to row.
\item To satisfy $\exists y \, \forall x \, P(x, y)$, an entire single column must be completely filled.
\end{itemize}

This is why $\exists y \, \forall x \, P(x, y)$ is a much stronger assertion than $\forall x \, \exists y \, P(x, y)$. In fact, the implication:
\[
\bigl(\exists y \, \forall x \, P(x, y)\bigr) \implies \bigl(\forall x \, \exists y \, P(x, y)\bigr)
\]
is always universally valid (if you have an unbroken column, every row clearly contains a point in that column).

Why is this implication guaranteed to hold? We can verify it directly using the variable roles introduced in Chapter~3:
\begin{enumerate}
\item \textbf{Assume the premise:} We assume that $\exists y \, \forall x \, P(x, y)$ is true. By the definition of existence, there is some specific witness $y^* \in D$ such that $P(x, y^*)$ holds for every $x \in D$.
\item \textbf{Set up the goal:} We need to prove $\forall x \, \exists y \, P(x, y)$. Following our game, let $x \in D$ be arbitrary (the Challenger hands us an element).
\item \textbf{Choose the response:} We must produce a winning response $y$. Which element do we pick? We choose $y = y^*$.
\item \textbf{Verify:} Does $P(x, y^*)$ hold? Yes, because $P(x, y^*)$ holds for \emph{all} elements in the domain, including the arbitrary $x$ handed to us.
\end{enumerate}
The reverse implication, however, is not generally true: having winning points in different rows does not mean they align in a single column.

\subsection{How the game dictates proof structure}

Quantifiers directly determine the logical skeleton of a proof (Figure~\ref{04:fig-proof-skeletons}):

\begin{figure}[htbp]
\centering
\begin{tikzpicture}[
    font=\small,
    proofcard/.style={draw=#1!70!black, fill=#1!4, rounded corners=4pt, thick, text width=6.6cm, inner sep=8pt}
]

% Left Card: Pointwise
\node[proofcard=targetgreen, anchor=north west] (cardA) at (0, 0) {
    \textbf{\color{targetgreen!90!black}Proving $\forall x \in D, \, \exists y \in D, \, P(x, y)$}\\[0.2em]
    \colorbox{targetgreen!80}{\color{white}\scriptsize\bfseries Pattern: Pointwise Response}\\[0.6em]
    \footnotesize
    \textbf{Round 1: Adversary moves ($\forall x$)}\\
    \emph{Proof line:} ``Let $x \in D$ be arbitrary.''\\
    \color{gray!80!black}Scope opens. $x$ is fixed but unknown; you cannot make any special assumptions about its value.\\[0.6em]
    \color{black}
    \textbf{Round 2: You respond ($\exists y$)}\\
    \emph{Proof line:} ``Choose $y = f(x)$.''\\
    \color{gray!80!black}You engineer a witness tailored to $x$. Your formula for $y$ is legally permitted to depend on $x$.\\[0.6em]
    \color{black}
    \textbf{Resolution: Verify the condition}\\
    \emph{Proof line:} ``Then $P(x, f(x))$ holds\dots''\\
    \color{gray!80!black}Show that this pair satisfies the predicate. Scope closes.
};

% Right Card: Uniform
\node[proofcard=sourceblue, anchor=north east] (cardB) at (14.6, 0) {
    \textbf{\color{sourceblue!90!black}Proving $\exists y \in D, \, \forall x \in D, \, P(x, y)$}\\[0.2em]
    \colorbox{sourceblue!80}{\color{white}\scriptsize\bfseries Pattern: Uniform Master Key}\\[0.6em]
    \footnotesize
    \textbf{Round 1: You move first ($\exists y$)}\\
    \emph{Proof line:} ``Choose $y = y^*$.''\\
    \color{gray!80!black}You must lock in a single candidate $y^*$ immediately. Your choice \emph{cannot} depend on $x$!\\[0.6em]
    \color{black}
    \textbf{Round 2: Adversary tests your key ($\forall x$)}\\
    \emph{Proof line:} ``Let $x \in D$ be arbitrary.''\\
    \color{gray!80!black}Scope opens. The adversary picks any element $x \in D$ in an effort to break your $y^*$.\\[0.6em]
    \color{black}
    \textbf{Resolution: Verify the condition}\\
    \emph{Proof line:} ``Then $P(x, y^*)$ holds\dots''\\
    \color{gray!80!black}Show your single key withstands every challenge $x$. Scope closes.
};

\end{tikzpicture}
\caption{Proof skeletons determined by quantifier order. In a $\forall x \, \exists y$ statement (left), the universal variable is introduced first and the existential witness may depend on $x$. In an $\exists y \, \forall x$ statement (right), the prover must produce a single witness $y^*$ that works for every $x$.}
\label{04:fig-proof-skeletons}
\end{figure}

\begin{recoverybox}
\textbf{Mental check:} In a proof of $\forall x \in \R, \, \exists y \in \R, \, P(x, y)$, can your definition of $y$ mention $x$? Can it mention a variable $z$ that has never been introduced?
\emph{Answer:} Yes to $x$, because $x$ was introduced in an outer scope before $y$ was chosen. No to $z$, because $z$ does not exist in the current scope. Variable scoping in proofs is just as strict as in computer programs.
\end{recoverybox}

\subsection{When can quantifiers swap? (Homogeneous vs.\ heterogeneous)}

Does quantifier order always matter? Not quite. The order matters only when you \textbf{switch players}:
\begin{itemize}
\item \textbf{Homogeneous quantifiers commute:}
\[
\forall x \, \forall y \, P(x, y) \iff \forall y \, \forall x \, P(x, y), \qquad\qquad \exists x \, \exists y \, P(x, y) \iff \exists y \, \exists x \, P(x, y).
\]
When both quantifiers are $\forall$, the Challenger chooses both elements. It makes no difference whether they pick $x$ first or $y$ first; their goal is simply to find any pair $(x, y)$ that breaks the claim. We often abbreviate this as $\forall x, y \in D$. Similarly, when both quantifiers are $\exists$, you get to pick both witnesses together.
\item \textbf{Heterogeneous quantifiers do not commute:}
\[
\forall x \, \exists y \, P(x, y) \not\iff \exists y \, \forall x \, P(x, y).
\]
As soon as $\forall$ and $\exists$ alternate, the game becomes asymmetric. Whoever moves second enjoys an information advantage: they get to inspect the first player's move before committing to their own.
\end{itemize}

\subsection{Everyday, software, and engineering analogies}

The distinction between $\forall x \, \exists y$ and $\exists y \, \forall x$ appears throughout technology and daily life.

\textbf{1. An audio amplifier specification:}
\begin{quote}
\textbf{Claim A}: For every frequency $f \in [20\,\text{Hz}, 20\,\text{kHz}]$, there exists an equalization filter setting that cancels distortion.
\end{quote}
This is a standard feature. The user turns a dial tailored to the specific problematic frequency. The setting responds to the frequency: $\text{setting} = g(f)$.

Now compare:
\begin{quote}
\textbf{Claim B}: There exists a single equalization filter setting that cancels distortion at every frequency $f \in [20\,\text{Hz}, 20\,\text{kHz}]$.
\end{quote}
This is a far more demanding engineering requirement. It demands a single static circuit that simultaneously eliminates all distortion across the whole audio spectrum.

\textbf{2. Software algorithms:}
\begin{itemize}
\item \emph{Pointwise:} ``For every input list $L$, there exists a Python script that sorts $L$ in $O(n)$ time.'' This is trivially true: if $L$ is already sorted, a simple verification script runs in linear time; if $L$ is in reverse order, a reverse-reading script works. You pick a custom script after inspecting the list!
\item \emph{Uniform:} ``There exists a comparison-based Python script that sorts every input list $L$ in $O(n)$ time.'' This is mathematically impossible: the famous comparison-sorting lower bound proves that any single comparison sort requires $\Omega(n \log n)$ operations in the worst case.
\end{itemize}

\textbf{3. Human language ambiguities:}
In everyday conversation, English often blurs quantifier order. If someone says ``Everyone loves someone,'' do they mean that each person has their own favorite companion ($\forall p \, \exists s$), or that there is a single beloved celebrity whom everyone admires ($\exists s \, \forall p$)? Mathematical notation exists precisely to prevent such ambiguities from corrupting our reasoning.

\begin{focusbox}
Whenever you read a definition with nested quantifiers, look at which variable appears first. A variable introduced later may depend on variables introduced earlier; a variable introduced earlier must be chosen independently of those that follow.
\end{focusbox}

\subsection{Three nested quantifiers and the bridge to negation}

The two-player game scales naturally to three, four, or any number of quantifiers: the players simply alternate turns.

Consider the definition of a \textbf{bounded function} $f \colon \R \to \R$ (a function whose values remain within a bounded interval):
\[
\exists M \in \R \text{ such that } \forall x \in \R, \, |f(x)| \le M.
\]
Here, you (the Prover) must move first by proposing a bound $M$. The Challenger then inspects your $M$ and tests whether any $x \in \R$ violates the bound. For $f(x) = \sin(x)$, you succeed by choosing $M = 1$. For $f(x) = x^2$, whatever $M$ you pick, the Challenger counters with $x = \sqrt{M} + 1$, so the function is unbounded.

Or consider the formal definition of the \textbf{limit of a sequence} $a_n \to L$:
\[
\forall \eps > 0, \, \exists N \in \N, \, \forall n \ge N, \, |a_n - L| < \eps.
\]
This corresponds to three rounds:
\begin{enumerate}
\item \textbf{Round 1 (Challenger, $\forall \eps > 0$):} Demands a tolerance, such as $\eps = 0.001$.
\item \textbf{Round 2 (Prover, $\exists N \in \N$):} Responds with a cutoff index $N$ computed from $\eps$.
\item \textbf{Round 3 (Challenger, $\forall n \ge N$):} Tests any subsequent term $n \ge N$ to see if $|a_n - L| < \eps$.
\end{enumerate}
If you have a rule for choosing $N$ that wins against every possible $\eps > 0$, the sequence converges.

We can visualize this three-round tournament as a sequence trapped within an epsilon band (Figure~\ref{04:fig-limit-game}).

\begin{figure}[htbp]
\centering
\begin{tikzpicture}[
    font=\small
]

% Epsilon band shading
\fill[targetgreen!10] (0.5, 1.4) rectangle (11.0, 2.6);

% Horizontal band boundaries
\draw[targetgreen!70!black, dashed, thick] (0.5, 2.6) -- (11.0, 2.6) node[right=2pt, font=\scriptsize, text=targetgreen!80!black] {$L + \eps$};
\draw[targetgreen!70!black, dashed, thick] (0.5, 1.4) -- (11.0, 1.4) node[right=2pt, font=\scriptsize, text=targetgreen!80!black] {$L - \eps$};

% Target limit L line
\draw[gray!50, dotted, thick] (1.5, 2.0) -- (11.0, 2.0) node[right=2pt, font=\scriptsize, text=gray!80] {$L$};

% Axes
\draw[->, thick, gray!70] (0.5, 0.4) -- (11.6, 0.4) node[right=3pt, font=\scriptsize] {$n$};
\draw[->, thick, gray!70] (0.5, 0.4) -- (0.5, 4.6) node[above=3pt, font=\scriptsize] {$a_n$};

% Epsilon width indicator
\draw[<->, thick, quotientorange] (1.1, 1.4) -- (1.1, 2.6);
\node[left=2pt, font=\tiny\bfseries, text=quotientorange!90!black] at (1.1, 2.0) {$2\eps$};

% Cutoff N vertical line
\draw[very thick, sourceblue, dashed] (6.0, 0.4) -- (6.0, 4.0);
\node[below=2pt, font=\scriptsize\bfseries, text=sourceblue!90!black] at (6.0, 0.4) {$N$};

% Round 1 callout
\node[anchor=south west, font=\scriptsize\bfseries, text=quotientorange!95!black] at (0.6, 4.2) {Round 1 ($\forall \eps$): Tolerance band $L \pm \eps$};

% Round 2 callout
\node[anchor=south west, font=\scriptsize\bfseries, text=sourceblue!90!black] at (6.2, 4.2) {Round 2 ($\exists N$): Cutoff index $N$};

% Round 3 callout
\node[anchor=north west, font=\scriptsize\bfseries, text=targetgreen!80!black, align=left] at (6.8, 1.25) {Round 3 ($\forall n \ge N$):\\All terms trapped inside!};

% Early terms (outside or wandering)
\fill[gray!50] (1.5, 3.6) circle (2.2pt) node[above=1pt, font=\tiny, text=gray!70] {$a_1$};
\fill[gray!50] (2.4, 0.8) circle (2.2pt) node[below=1pt, font=\tiny, text=gray!70] {$a_2$};
\fill[gray!50] (3.3, 3.1) circle (2.2pt) node[above=1pt, font=\tiny, text=gray!70] {$a_3$};
\fill[gray!50] (4.2, 1.1) circle (2.2pt) node[below=1pt, font=\tiny, text=gray!70] {$a_4$};
\fill[gray!50] (5.1, 2.8) circle (2.2pt) node[above=1pt, font=\tiny, text=gray!70] {$a_5$};

% Terms from N onward (safely trapped)
\fill[targetgreen!80!black] (6.0, 2.2) circle (2.4pt);
\node[right=3pt, font=\tiny\bfseries, text=targetgreen!80!black] at (6.0, 2.2) {$a_N$};
\fill[targetgreen!80!black] (7.0, 1.75) circle (2.2pt);
\fill[targetgreen!80!black] (8.0, 2.20) circle (2.2pt);
\fill[targetgreen!80!black] (9.0, 1.88) circle (2.2pt);
\fill[targetgreen!80!black] (10.0, 2.08) circle (2.2pt);

\end{tikzpicture}
\caption{The three-round quantifier game for the limit of a sequence ($\forall \eps > 0, \exists N \in \N, \forall n \ge N, |a_n - L| < \eps$). In Round~1, the Challenger narrows the tolerance corridor to $L \pm \eps$. In Round~2, the Prover inspects $\eps$ and chooses a cutoff index $N$. In Round~3, the Challenger tests any $n \ge N$: all terms from $N$ onward satisfy the condition $|a_n - L| < \eps$.}
\label{04:fig-limit-game}
\end{figure}

What if a quantified claim is false? How do you mathematically disagree with an assertion containing multiple quantifiers?

If an adversary claims that $\forall x \in D, \, \exists y \in D, \, P(x, y)$, refuting them requires you to step into the shoes of the Challenger. You must establish that \emph{there exists an element $x \in D$} for which \emph{every possible response $y$} fails:
\[
\neg \bigl(\forall x \in D, \, \exists y \in D, \, P(x, y)\bigr) \iff \exists x \in D, \, \forall y \in D, \, \neg P(x, y).
\]
Notice the systematic inversion: every $\forall$ transforms into an $\exists$, every $\exists$ transforms into a $\forall$, and the condition is negated. Negation exchanges the roles of the universal and existential quantifiers.

In Chapter~5 (\emph{Negation: How to Disagree Precisely}), we examine the systematic rules for negating quantified statements and constructing counterexamples.

\subsection{Exercises}

\subsubsection*{Part A: Warm-Up Mental Checks}
State whether each statement is true or false, and justify your answer in a single sentence:
\begin{enumerate}
\item $\forall x \in \R, \, \exists y \in \R \text{ such that } x + y = 10$.
\item $\exists y \in \R \text{ such that } \forall x \in \R, \, x + y = 10$.
\item $\forall x \in \Z, \, \exists y \in \Z \text{ such that } x - y \text{ is an even integer}$.
\end{enumerate}

\subsubsection*{Part B: Routine Game Practice}
Determine whether each statement is true or false. Provide a winning strategy for the Prover, or an adversary's winning counter-move:
\begin{enumerate}[resume]
\item $\forall x \in \R, \, \exists y \in \R \text{ such that } x < y$.
\item $\exists y \in \R \text{ such that } \forall x \in \R, \, x < y$.
\item $\forall x \in \R, \, \exists y \in \R \text{ such that } xy = 1$.
\item $\forall x \in (\R \setminus \{0\}), \, \exists y \in \R \text{ such that } xy = 1$.
\item $\exists y \in \R \text{ such that } \forall x \in \R, \, xy = 0$.
\item $\forall x \in \N, \, \exists y \in \N \text{ such that } y < x$ (where $\N = \{1, 2, 3, \ldots\}$).
\end{enumerate}

\subsubsection*{Part C: Proof Skeletons}
For each claim below, write only the opening declaration (introducing variables and their roles) and the final verification line of a direct proof. Do not carry out the calculations:
\begin{enumerate}[resume]
\item $\forall x \in \R, \, \exists y \in \R \text{ such that } y^3 = x$.
\item $\exists c \in \R \text{ such that } \forall x \in \R, \, x^2 + c \ge 0$.
\end{enumerate}

\subsubsection*{Part D: Deeper Challenges}
\begin{enumerate}[resume]
\item Write out a rigorous four-step proof that for any non-empty domain $D$ and any predicate $P(x, y)$, the implication $\bigl(\exists y \in D \, \forall x \in D \, P(x, y)\bigr) \implies \bigl(\forall x \in D \, \exists y \in D \, P(x, y)\bigr)$ is valid. Explicitly state the role (arbitrary vs.\ chosen) of each variable you introduce.
\item Consider the three-quantifier sentence on the real numbers:
\[
\forall x \in \R, \, \exists y \in \R, \, \forall z \in \R, \, x + y + z > 0.
\]
Is this statement true or false? Play the three-round game step by step to prove your conclusion.
\end{enumerate}

\begin{walkbox}
Think of the definition of the limit of a sequence: ``For every $\eps > 0$, there exists an index $N$ such that\ldots'' Who moves first in this game? Does $N$ depend on $\eps$, or does $\eps$ depend on $N$? What would it mean if the quantifiers were reversed: ``There exists an index $N$ such that for every $\eps > 0$\ldots''? Could any sequence except an eventually constant one satisfy that reversed condition?
\end{walkbox}

\newpage
